\subsubsection{Unison \& Polyphony}
\label{sec:unison}

The unison effect generates multiple voices of the same sound, each slightly detuned,
therefore creating a richer and thicker sound.  Referring from other oscillator-based synthesizers,
for e.g. Xfer Serum, the equation for creating unison voices can be modelled like this, 
with $x[n]$ as the input oscillator of any shape and frequency:

\begin{equation}
    y[n] = A_c x[n] + A_s \sum_{k=1}^{N} \left( x\left[ n f(d, k) - \phi_k \right]
            +  x\left[ n f(-d, k) - \phi_k \right] \right)
\end{equation}

Where $A_c$ is the amplitude of the centre voice and $A_s$ is the amplitude of the 
side voices. $f(k)$ is the detune factor for the frequency, $d$ is the detune amount (how 
far the frequencies are detuned from each other), and $\phi$ is the phase shift as described in 
\eqref{eq:unison_phase}.

The detune factor is described using the following equation:

\begin{equation}\label{eq:cent_frequency}
f(d, k) = 2^\frac{dk}{1200}
\end{equation}

Where $k$ is the cent value, which defines the logarithmic unit in music. Szabo's paper makes a discussion
that the detuned frequency changes non-linearly with the interval length, $d$ \cite{szabo_supersaw}. 
Since there is no documentation published online, using this equation for detuning \eqref{eq:cent_frequency}
would be valid (\Cref{prompt:unison_detune}) \cite{openai_chatgpt}. Also adding a cent value add a slight, insignificant
change in the frequency without altering the perceived pitch, allowing multiple voices to remain sounding natural.

Regarding the design of blending of the voices, the empirical analysis of the Roland JP-8000
in Szabo's paper suggests that the centre oscillator must decrease proportionally along
with the increase in side oscillator counterparts \cite{szabo_supersaw}, to ensure proper
distribution of energy, when the voices are spreading.

After the oscillator wave is splitted into voices, an equation using energy conservation is derived,
given that the power, $P \propto A^2$ (\Cref{prompt:unison_gainblend}) \cite{anthropicclaude}, 

\begin{equation}\label{eq:unison_energyconserv}
(A_c)^2 + (N-1)(A_s)^2 = 1
\end{equation}

where $N$ is the total number of voices.

The blending factor, $b$ where $0 \leq b \leq 1$ is applied to the power space, which makes
$P_c = (A_c)^2 = 1 - b$.

$$
1 - b + N(A_s)^2 = 1
$$

\begin{equation}\label{eq:unison_equalpower}
\therefore A_s = \sqrt{\frac{b}{N-1}}, \qquad A_c = \sqrt{1 - b}
\end{equation}

This equation is only applicable for odd-numbered voices. For even numbered voices with two
centre voices, 

\begin{equation}\label{eq:unison_energyconserveven}
2(A_c)^2 + (N-1)(A_s)^2 = 1
\end{equation}

After expressing $P_c = 2(A_c)^2 = 1 - b$, the following is obtained:
\begin{equation}\label{eq:unison_equalpower_even}
\therefore A_s = \sqrt{\frac{b}{N-2}}, \qquad A_c = \sqrt{ \frac{1 - b}{2} }
\end{equation}

These equations can also be derived using the wave superposition properties that the sum of 
all the amplitudes of waves must equal the amplitude of the input.

\begin{equation}\label{eq:unison_wavesuperpos}
A_c + NA_s = 1
\end{equation}

This time, $b$ is applied in the amplitude domain, where $A_c = 1 - b$.
$$
1 - b + NA_s = 1
$$

\begin{equation}\label{eq:unison_linearmode}
\therefore A_s = \frac{b}{N}, \qquad A_c = 1 - b
\end{equation}

The linear blend equation \eqref{eq:unison_linearmode} makes it more intuitive and easy to implement, but unlike the equal-power
counterpart \eqref{eq:unison_equalpower}, it doesn't preserve total signal power across the blending range, and leads to 
perceived loudness \cite{anthropicclaude}. Even though Szabo's Roland instrument demonstrated exponential blend, it is not power-preserving
and that the side voices stay softened. Therefore, with that being mentioned it is safe to conclude that the equal-power blend provides
the most robust solution, even though it can be switched in the next version.

On the other hand, the phase change, $\phi$ is assigned randomly as described in Szabo's paper, and is defined by the following:

\begin{equation}\label{eq:unison_phase}
    \phi_i \sim \mathcal{U}\!\left(0,\,2\pi\lambda\right), \qquad 0 < \lambda \leq 1
\end{equation}

The adjustable scalar parameter $\lambda$ controls the boundary within which the values can be assigned. $\lambda = 0$ provides consistent
phases, whereas $\lambda = 1$ gives maximum variation. The greater the variation, the more natural and organic the synth sounds \cite{szabo_supersaw}.

However, to test the perceptual quality of the sound, a sawtooth wave is passed into the system, which then duplicates into 3 voices,
detuned for about 10 cents from each other, and a gain ratio of 0.7. Starting with a low blending factor $b$, it produced a richer, thicker sound 
but the beats produced clipping artifacts, which is unexpected. The following table shows how the sound quality is
optimized through adjusting specific parameters.

\begin{table}[h]
\centering
\caption{Observations during refinement of the unison implementation.}
\begin{tabular}{|l|p{2.5cm}|>{\centering\arraybackslash}p{4cm}|p{5.5cm}|}
\hline
\textbf{Parameter} &
\textbf{Adjustment} &
\textbf{Clipping Noticeable?} &
\textbf{Observation} \\
\hline

Phase Shift &
Random &
Sometimes &
Randomization altered the interference pattern between oscillators, resulting in different beating characteristics. \\
\hline

Blend Factor &
Increased to 1.0 &
No &
Redistribution of the gain among the centre and side voices reduces the intensity of the beating effect. \\
\hline

Master Gain &
Reduced from 0.7 to 0.4 &
No &
Reducing the master gain prevented clipping, thereby eliminating audible distortions. \\
\hline

\end{tabular}
\label{tab:unison_observations}
\end{table}

Therfore, it can be concluded that the unison produces expected result, if the gain is reduced
to prevent any clippings driven by beating. But the quality can further be improved either
by rearranging the shifts in another way or by spreading the voices throughout stereo output, 
since it's currently working in mono, which can reduce wobbly beating effect. Code listing is included
under \Cref{app:unison_code}.
